\section{Genus two: compact polyhedra and a marked-torus model}
\label{sec:genus2markedtorus}

The hyperelliptic quotient makes genus two a concrete testing ground
for the counting converse.  There are two useful reductions.  The
hyperbolic quotient gives a compact polyhedral model for every
conformal base.  On the nonsquare stratum with two double zeros,
a second, flat quotient identifies the entire stratum flow space, up to
finite markings, with the admissible invariant open subset of an affine
homogeneous space of marked flat tori. Its complement has measure zero.
The latter identification is specific to that stratum.

\subsection{Six branch points and the compact polyhedron}
\label{subsec:genus2compact}

Every genus-two Riemann surface $X$ is hyperelliptic.  Its quotient
map $\pi:X\to\mathbb P^1$ has six branch values: the
Riemann--Hurwitz formula reads $2=2(-2)+6$.  With compatible
markings the quotient induces a holomorphic Teichm\"uller isometry
\[
 \T_{2,0}\cong\T_{0,6};
\]
see, for example, \cite[Corollary~1.2(3)]{BenirschkeServan}.
Both sides have complex dimension three.  The four-pointed sphere
is the corresponding genus-one example.

The hyperelliptic involution is an isometry of the curvature-$-1$
metric on $X$.  The quotient metric is a hyperbolic sphere with
\emph{six cone angles $\pi$}, and its area is $2\pi$.
Alexandrov's hyperbolic realization theorem therefore gives a
unique compact convex polyhedron in $\HH^3$ whose intrinsic
boundary metric is this quotient metric, allowing the usual
degenerate case of a doubled convex polygon
\cite[Theorem~1.1]{FillastreIzmestievCusps}.
There are six vertices, all with intrinsic total angle $\pi$.
Conversely, take the double cover of such a cone sphere branched at
its six vertices.  Each cone angle becomes $2\pi$, so the lifted
metric is smooth, complete, and of curvature $-1$.  Riemann--Hurwitz
identifies its genus as two.  Thus this compact-polyhedral
realization is an exact model for the genus-two hyperbolic metric,
with the appropriate finite labeling and marking conventions.

This metric must be distinguished from the complete hyperbolic
metric on the sphere with the six branch values deleted.  The
latter has six cusps and area $8\pi$; it is the metric to which the
ideal-polyhedron realization applies.  Passing between them uses
uniformization within the same punctured conformal class.  It does
not identify their lengths, vertex data, or volume forms.
In particular, fixed intrinsic cone angles at compact vertices are
not automatically linear constraints on the dihedral angles.

\subsection{An explicit formula for every quadratic differential}
\label{subsec:genus2quadraticformula}

Choose a coordinate with all six branch values finite and write
\begin{equation}
 X:\quad y^2=B_6(z):=\prod_{i=1}^{6}(z-a_i).
 \label{eq:genus2curve}
\end{equation}
The forms $\dd z/y$ and $z\dd z/y$ form a basis of
$H^0(X,K_X)$.  Their three symmetric products are linearly
independent and, by Riemann--Roch, span $H^0(X,K_X^2)$.
Hence every holomorphic quadratic differential is
\begin{equation}
 q_X=P_2(z)\frac{\dd z^2}{y^2}
     =\pi^*q_0,
 \qquad
 q_0=\frac{P_2(z)}{B_6(z)}\dd z^2,
 \qquad \deg P_2\leq2.
 \label{eq:genus2qformula}
\end{equation}
In particular, every such differential is invariant under the
hyperelliptic involution.  At a branch point the order transforms by
\begin{equation}
 \operatorname{ord}_p(\pi^*q_0)
 =2\operatorname{ord}_{\pi(p)}(q_0)+2.
 \label{eq:genus2order}
\end{equation}
For a generic polynomial with two distinct roots away from the
branch values, this gives
\[
 \mathcal Q_2(1^4)\longleftrightarrow
 \mathcal Q_0(1^2,-1^6),
\]
with finite choices of labels understood.  This is the genus-two
case of the hyperelliptic covering construction in
\cite[Section~1.1.3]{LanneauComponents}.
The descended flat metric has six angles $\pi$ at its poles and
two angles $3\pi$ at its zeros.  It is different from the
hyperbolic cone metric of the preceding subsection.

Pullback multiplies flat area by two.  Thus the correspondence
between unit-area differentials on the six-pointed sphere and on
$X$ is
\[
 q_0\longmapsto\tfrac12\pi^*q_0.
\]
It commutes with $g_t=\diag(e^t,e^{-t})$, with no change of the time
parameter.  Period-lattice covolumes, finite labelings, and orbifold
stabilizers must still be tracked when comparing unnormalized
Masur--Veech volumes; the area factor alone is not a formula for
those volumes.

\subsection{The nonsquare stratum with two double zeros}
\label{subsec:genus2twodouble}

Write $\mathcal Q^{\mathrm{ns}}_2(2,2)$ for quadratic
differentials with two double zeros which are not global squares
of Abelian differentials.  The superscript is essential: the
squares of differentials in $\mathcal H_2(1,1)$ have the same
quadratic zero orders but form a different locus.

Formula~\eqref{eq:genus2qformula} gives an elementary description.
The differential is a global square exactly when $P_2$ is the square
of a polynomial of degree at most one: this follows by writing every
Abelian differential as $(az+b)\,\dd z/y$. In homogeneous degree-two
notation this includes a double root at infinity. A nonzero constant
$P_2$ is therefore a square, while a degree-one $P_2$ has a simple zero
at each of the two unbranched points at infinity and cannot lie in
$\mathcal Q_2(2,2)$. For degree two, a simple root at a branch value
gives one double zero by~\eqref{eq:genus2order}, and a simple root
elsewhere gives two simple zeros. Thus, to obtain two double zeros
without a square, the roots must be two distinct branch values, say
$b_1,b_2$.
After relabeling the remaining four branch values,
\begin{equation}
 \begin{split}
 X:&\quad y^2=B_4(z)(z-b_1)(z-b_2),
       \qquad B_4(z)=\prod_{i=1}^4(z-a_i),\\
 q_X&=c(z-b_1)(z-b_2)\frac{\dd z^2}{y^2}
      =\pi^*\left(\frac{c\,\dd z^2}{B_4(z)}\right),
      \qquad c\in\CC^*.
 \end{split}
 \label{eq:genus2cancel}
\end{equation}
The six values $a_1,\ldots,a_4,b_1,b_2$ are distinct.
Downstairs $q_0=c\,\dd z^2/B_4(z)$ has four simple poles and is
regular at the two additionally marked points $b_1,b_2$:
\begin{equation}
 \mathcal Q^{\mathrm{ns}}_2(2,2)
 \longleftrightarrow \mathcal Q_0(-1^4,0,0).
 \label{eq:genus2pillowmark}
\end{equation}
The regular marks cannot be discarded.  They are branch values of
$X\to\mathbb P^1$, and their positions supply four real
parameters.

For each fixed $X$, choosing the unordered pair of branch values
$b_1,b_2$ gives $\binom62=15$ complex lines of nonsquare
$(2,2)$ differentials in the three-dimensional cotangent space.
Area normalization leaves a phase on each line.  Consequently this
stratum meets every conformal base, but its unit-area locus has
real dimension seven, compared with eleven for the full principal
unit cotangent locus.  Its Masur--Veech measure is an intrinsic
stratum measure, not a conditional probability obtained by
restricting the principal measure to a set of measure zero.

\subsection{The pillowcase cover and two marked torus points}
\label{subsec:genus2ellipticcover}

The orientation double cover of the four-pole differential $q_0$
is the elliptic curve
\begin{equation}
 E:\quad w^2=B_4(z),
 \qquad \omega=\sqrt c\,\frac{\dd z}{w},
 \qquad \rho^*q_0=\omega^2.
 \label{eq:genus2toruscover}
\end{equation}
Choose one branch point of $E\to\mathbb P^1$ as the origin.
The involution is then $v\mapsto-v$, and its four fixed points are
$E[2]$.  The two marked regular points lift to unordered pairs
\[
 \rho^{-1}(b_1)=\{p,-p\},\qquad
 \rho^{-1}(b_2)=\{q,-q\}.
\]
Their admissibility conditions are exactly
\begin{equation}
 p,q\notin E[2],\qquad p\ne q,\qquad p\ne -q.
 \label{eq:genus2torusadmissible}
\end{equation}
Conversely, a flat torus $(E,\omega)$ with an origin and such two
points gives the sphere $E/\{\pm1\}$ with its four poles and two
regular marks.  The unique connected double cover of this sphere
branched at all six marked values is $X$, and pullback of $q_0$
gives~\eqref{eq:genus2cancel}.
The cover is connected because its monodromy is the nontrivial
homomorphism to $\ZZ/2\ZZ$ on every small loop around a branch
value; six such loops satisfy the sphere relation.  It is unique
because those loops generate the fundamental group of the
punctured sphere.

The resulting $q_X$ is nonsquare.  Indeed, its being square would
force the orientation cover $E\to\mathbb P^1$ to factor through
$X\to\mathbb P^1$.  Both covers have degree two, but their branch
sets differ: the former is branched at the four $a_i$, the latter
also at $b_1,b_2$.
Equivalently, the square root of $B_4$ is not in the function field
$\CC(z,\sqrt{B_4(z)(z-b_1)(z-b_2)})$.

Both covering maps to the sphere have degree two, so
\begin{equation}
 \operatorname{Area}(E,\omega)
 =2\operatorname{Area}(q_0)
 =\operatorname{Area}(X,q_X).
 \label{eq:genus2equalareas}
\end{equation}
Thus the unit-area marked torus corresponds directly to the
unit-area genus-two differential, with no further area scaling.
The entire construction commutes with the $SL_2(\RR)$ action on
flat charts.  In particular its time parameter is exactly the
Teichm\"uller time used above.


% BEGIN INLINE FIGURE: 06_double_covers.tex
\begin{figure}[htbp]
\centering
\begin{tikzpicture}[font=\small,>=Latex,line cap=round,line join=round]
  \node[font=\bfseries\small] at (-3.5,3.75) {Genus-two surface $X$};
  \begin{scope}[shift={(-3.5,2.15)}]
    \filldraw[fill=mpblue!7,draw=mpblue,line width=.9pt]
      (-2.2,0) .. controls (-2.2,1) and (-1.4,1.25) .. (-.75,.85)
      .. controls (-.3,.55) and (.3,.55) .. (.75,.85)
      .. controls (1.4,1.25) and (2.2,1) .. (2.2,0)
      .. controls (2.2,-1) and (1.4,-1.2) .. (.75,-.8)
      .. controls (.3,-.5) and (-.3,-.5) .. (-.75,-.8)
      .. controls (-1.4,-1.2) and (-2.2,-1) .. (-2.2,0);
    \filldraw[fill=white,draw=mpblue] (-1.2,.05) ellipse (.58 and .23);
    \filldraw[fill=white,draw=mpblue] (1.2,.05) ellipse (.58 and .23);
    \draw[mpblue!55] (-1.78,.05) .. controls (-1.5,-.6) and (-.9,-.6) .. (-.62,.05);
    \draw[mpblue!55] (.62,.05) .. controls (.9,-.6) and (1.5,-.6) .. (1.78,.05);
  \end{scope}
  \node[align=center,font=\scriptsize,text width=6cm] at (-3.5,.65)
    {$y^2=B_4(z)(z-b_1)(z-b_2)$\\$q_X=\pi^*q_0\in\mathcal Q_2^{\mathrm{ns}}(2,2)$};
  \node[font=\bfseries\small] at (3.5,3.75) {Marked flat torus $E$};
  \begin{scope}[shift={(1.75,1.05)},x={(3cm,0cm)},y={(.6cm,2.15cm)}]
    \filldraw[fill=mpteal!5,draw=mpteal,line width=.8pt] (0,0)--(1,0)--(1,1)--(0,1)--cycle;
    \draw[mpteal!25] (.5,0)--(.5,1) (0,.5)--(1,.5);
    \draw[->,mpteal] (.35,0)--(.65,0);\draw[->,mpteal] (.35,1)--(.65,1);
    \draw[->,mpteal] (0,.35)--(0,.65);\draw[->,mpteal] (1,.35)--(1,.65);
    \foreach \p in {(0,0),(.5,0),(0,.5),(.5,.5)}\fill[mpblue] \p circle (1.8pt);
    \fill[mpochre] (.2,.3) circle (2pt) node[left=2pt,font=\scriptsize] {$p$};
    \fill[mpochre] (.8,.7) circle (2pt) node[right=2pt,font=\scriptsize] {$-p$};
    \fill[mpred] (.25,.75) circle (2pt) node[left=2pt,font=\scriptsize] {$q$};
    \fill[mpred] (.75,.25) circle (2pt) node[right=2pt,font=\scriptsize] {$-q$};
  \end{scope}
  \node[align=center,font=\scriptsize,text width=5.5cm] at (3.5,.65)
    {$w^2=B_4(z),\qquad \omega=\sqrt c\,dz/w$\\$\rho^*q_0=\omega^2$};
  \draw[->,mpgray,line width=.8pt] (-2.9,.15)--(-1.05,-1.0)
    node[midway,above left,font=\scriptsize,align=center] {degree two\\six branch values};
  \draw[->,mpgray,line width=.8pt] (2.9,.15)--(1.05,-1.0)
    node[midway,above right,font=\scriptsize,align=center] {degree two\\four branch values};
  \begin{scope}[shift={(0,-2.25)}]
    \filldraw[fill=mpblue!3,draw=mpgray!70] (0,0) circle (1.35);
    \draw[mpgray!25,dashed] (-1.35,0) arc (180:360:1.35 and .32);
    \draw[mpgray!35] (-1.35,0) arc (180:0:1.35 and .32);
    \fill[mpblue] (-.92,.51) circle (2pt) node[left=2pt,font=\scriptsize] {$a_1$};
    \fill[mpblue] (-.63,-.72) circle (2pt) node[below left=1pt,font=\scriptsize] {$a_2$};
    \fill[mpblue] (.68,-.7) circle (2pt) node[below right=1pt,font=\scriptsize] {$a_3$};
    \fill[mpblue] (.93,.47) circle (2pt) node[right=1pt,font=\scriptsize] {$a_4$};
    \fill[mpochre] (-.23,.43) circle (2.3pt) node[above=2pt,font=\scriptsize] {$b_1$};
    \fill[mpred] (.25,-.3) circle (2.3pt) node[right=2pt,font=\scriptsize] {$b_2$};
  \end{scope}
  \node[align=center] at (0,-4.3)
    {$q_0=\displaystyle\frac{c\,dz^2}{B_4(z)}$ on $\mathbb P^1$\\[4pt]
     \textcolor{mpblue}{four simple poles $a_i$},
     \textcolor{mpochre}{two regular marks $b_1,b_2$}};
\end{tikzpicture}
\caption[The genus-two and marked-torus double covers]{The two double covers used on the nonsquare genus-two stratum
$\mathcal Q_2^{\mathrm{ns}}(2,2)$. The cover $X\to\mathbb P^1$ branches
at all six marked values; the orientation cover $E\to\mathbb P^1$
branches only at the four poles. The two regular marks lift to the pairs
$\{p,-p\}$ and $\{q,-q\}$. Blue points in the torus diagram represent
the four classes in $E[2]$, with boundary identifications understood.
This is a diagram of flat differentials, distinct from the hyperbolic
six-cone-point quotient metric.}
\label{fig:mono-double-covers}
\end{figure}
% END INLINE FIGURE

\subsection{A homogeneous finite cover}
\label{subsec:genus2homogeneous}

Choose a full symplectic level-two structure on the torus and
choose an ordering and one lift from each of the two pairs
$\{\pm p\},\{\pm q\}$.  Together with the origin, level two
labels the four two-torsion points and hence the four poles.  Set
\[
 \Gamma_2=\ker\bigl(SL_2(\ZZ)\to SL_2(\ZZ/2\ZZ)\bigr),
 \qquad
 G_2=SL_2(\RR)\ltimes(\RR^2\oplus\RR^2),
\]
where $SL_2(\RR)$ acts in the standard way on both vector factors.
Our group law and right-quotient convention are
\begin{align}
 (A,u_1,u_2)(B,v_1,v_2)
   &=(AB,u_1+Av_1,u_2+Av_2),\\
 Y_2&=G_2/(\Gamma_2\ltimes(\ZZ^2\oplus\ZZ^2)).
 \label{eq:genus2Y2}
\end{align}
A representative $(A,u_1,u_2)$ describes the area-one torus
$\RR^2/A\ZZ^2$ with marked points $[u_1],[u_2]$.
Right multiplication by $(\gamma,n_1,n_2)$ changes the
representative to
\[
 (A\gamma,u_1+An_1,u_2+An_2),
\]
which is exactly change of the lattice basis preserving level two
and change of the two point lifts.  The presence of $-I$ in
$\Gamma_2$ causes no singularity in this homogeneous quotient:
right translation by a discrete subgroup on a Lie group is free.
Here $A$ and $-A$ give the same physical lattice and the same
level-two data.  A finer congruence level gives an additional finite cover. The spectral
argument below uses level two specifically; its conclusion about the
base spectrum is not automatically inherited by further covers.

Delete the locus where one of the points is two-torsion or where
$[u_1]=\pm[u_2]$, and call the resulting open set $Y_2^\circ$.
In each torus fiber the deleted set is a finite union of real
codimension-two subtori.  It has Haar measure zero and its
complement is connected.  Since the base $SL_2(\RR)/\Gamma_2$ is
connected, so is $Y_2^\circ$.
Forgetting the finite marking choices gives a finite equivariant
local parameterization of $\mathcal Q^{\mathrm{ns},1}_2(2,2)$.
The finite choices include the distinguished origin over a pole,
the level structure, the ordering of the two regular marks, and
their signs.  Their changes preserve Haar measure and commute
with the $SL_2(\RR)$ action; the stratum is their finite quotient.
Over the coarse-moduli locus without additional automorphisms
this is an ordinary finite covering.  If one retains the
ubiquitous genus-two hyperelliptic involution in the orbifold
stabilizers, its fixed order-two contribution must also be
included in unnormalized volume conventions.  No infinite
marking choice is being discarded in this finite-cover statement.

The flow is left multiplication:
\begin{equation}
 g_t[A,u_1,u_2]
   =[a_tA,a_tu_1,a_tu_2],
 \qquad a_t=\diag(e^t,e^{-t}).
 \label{eq:genus2homflow}
\end{equation}
The admissible set $Y_2^\circ$ is invariant under this action.
Although the coordinates $\xi_i=A^{-1}u_i\pmod{\ZZ^2}$ are
unchanged in a fixed lattice chart, a change $A\mapsto A\gamma$
changes them by $\xi_i\mapsto\gamma^{-1}\xi_i$.  Therefore they
are not globally fixed torus coordinates on the quotient.
This monodromy is essential in the later Fourier calculation.

\subsection{Haar measure and the period lattice}
\label{subsec:genus2haarmeasure}

The group $G_2$ is unimodular, and its Haar measure is
\[
 \dd A\,\dd u_1\,\dd u_2,
\]
where $\dd A$ is Haar measure on $SL_2(\RR)$.
In a local lattice chart, $u_i=A\xi_i$ and $\det A=1$, so its
quotient measure is the product of modular Haar measure and
Lebesgue probability on $(\RR^2/\ZZ^2)^2$.
The modular quotient has finite Haar volume and each torus fiber
has volume one, hence $Y_2$ has finite volume.

Before area normalization, local coordinates are the two complex
lattice periods and the two complex positions of the marked
points.  These are four complex coordinates, agreeing with
\[
 \dim_{\CC}\mathcal Q^{\mathrm{ns}}_2(2,2)=4.
\]
The covering construction transports these coordinates to
quadratic-differential period coordinates by a fixed linear map
on each marked chart.  This can be seen without a global
classification theorem.  The orientation cover of $q_X$ is the
fiber product
\[
 \widehat X:\quad
 w^2=B_4(z),\qquad t^2=(z-b_1)(z-b_2),
 \qquad y=wt,
 \qquad \widehat\omega=\sqrt c\,\frac{\dd z}{w}.
\]
Its map to $E$ is a double cover branched at
$\{p,-p,q,-q\}$.  Thus its relative periods are integral linear
combinations of torus periods and marked-point positions, after
paths and lifts have been fixed.  Transfer on absolute cycles and
paths between $p$ and $-p$, or between $q$ and $-q$, recover the
torus periods and twice the point positions, up to periods.
Consequently the inverse change of coordinates is rational
linear.  The integral period lattices are commensurable; factors
of two can occur, so the map need not be unimodular.  This is the
period-coordinate construction underlying the stratum measure
\cite[Sections~1.1.1--1.1.2]{LanneauComponents}.

It follows that the transported Masur--Veech measure is a constant
multiple of Lebesgue measure in these cone coordinates.  More
explicitly, write the torus area as $s>0$, put $A=\sqrt{s}\,B$
with $\det B=1$, and use the point coordinates $\xi_i=A^{-1}u_i$.
Up to a fixed positive constant the cone measure is
\[
 s^3\,\dd s\,\dd B\,\dd\xi_1\,\dd\xi_2.
\]
Restricting by area as in~\eqref{eq:genus2equalareas} therefore
makes the unit-area measure a constant multiple of Haar measure
on $Y_2$.  After
normalizing to probability and accounting for the finite covering,
the measures agree.  No numerical identity between unnormalized
volumes is asserted without a choice of lattice and labeling
conventions.  The measure-zero exclusions in $Y_2\setminus
Y_2^\circ$ do not affect this probability identification.

For observables compactly supported inside $Y_2^\circ$, extension
by zero gives smooth observables on the homogeneous space.  Thus
estimates proved there pass to the stratum after lifting the
observables and, if necessary, averaging over the finite markings.
For observables which grow near collisions, extension and norm
control require additional estimates; a null collision set alone
is not a bound on such observables.


\subsection{Fourier decomposition and temperedness}
\label{subsec:markedtorusconverse}

The affine homogeneous model admits a direct calculation which does
not use a general Teichm\"uller mixing theorem. Let
\[
 G=\operatorname{SL}_2(\RR),\qquad
 \mathsf A=G\ltimes\operatorname{Mat}_{2\times2}(\RR),\qquad
 \Lambda=\Gamma\ltimes\operatorname{Mat}_{2\times2}(\ZZ),
\]
where $\Gamma=\Gamma(2)$ initially. The group law is
\[
 (g,v)(g',v')=(gg',v+gv').
\]
Set $Y_2=\mathsf A/\Lambda$, with Haar probability measure $\mu$, and
let $a_t=\diag(e^t,e^{-t})$ act by left translation. A point is a
unit-area lattice together with two marked vectors modulo that
lattice. If one counts the origin as a marked point, this is a
torus with three ordered marked points. Coincidence loci have
measure zero and may be removed. Any application to a particular
quadratic stratum uses its separately established geometric
identification with this model or a finite quotient; an arbitrary
finite cover cannot be discarded without checking its base spectrum.

\subsubsection*{Fiber Fourier coefficients and their stabilizers}

Write $v=gu$, with $u\in\operatorname{Mat}_{2\times2}(\RR)/
\operatorname{Mat}_{2\times2}(\ZZ)$. Then
\[
 (g,u)\sim(g\gamma,\gamma^{-1}u),\qquad
 a_t(g,u)=(a_tg,u).
\]
For $F\in L^2(Y_2)$ expand, in the $L^2$ sense,
\begin{equation}
 F(g,u)=\sum_{M\in\operatorname{Mat}_{2\times2}(\ZZ)}
   F_M(g)e^{2\pi i\tr(M^Tu)}.
 \label{eq:markedtorusfourier}
\end{equation}
The equivariance relation is
\begin{equation}
 F_M(g\gamma)=F_{\gamma^{-T}M}(g).
 \label{eq:markedtorusequivariance}
\end{equation}
Consequently each $\Gamma$-orbit of frequency matrices contributes
a $G$-invariant summand isometric to
\begin{equation}
 L^2(G/\Gamma_M),\qquad
 \Gamma_M=\{\gamma\in\Gamma:\gamma^TM=M\}.
 \label{eq:markedtorusinduction}
\end{equation}
Indeed, Parseval followed by unfolding the orbit sum gives
\[
 \int_{G/\Gamma}
       \sum_{M'\in\Gamma\cdot M}|F_{M'}(g)|^2\,\dd g
   =\int_{G/\Gamma_M}|F_M(g)|^2\,\dd g.
\]

The stabilizers are elementary. If $M=0$, the summand is
$L^2(G/\Gamma)$. If $M$ has rank two, $\Gamma_M=\{1\}$.
If $M$ has rank one, write $M=vw^T$, with $v$ a primitive
integral column vector and $w\ne0$ an integral column vector.
Then $\Gamma_M$ fixes $v$ under the transpose action. It is a
subgroup conjugate to a discrete unipotent group, hence is
amenable. For example, if $v=(1,0)^T$ and $\Gamma=\Gamma(2)$,
\[
 \Gamma_M=
 \left\{\begin{pmatrix}1&0\\2n&1\end{pmatrix}:n\in\ZZ\right\}.
\]
No transitivity assertion about primitive vectors under $\Gamma(2)$
is required: one simply sums over all its frequency orbits. The same
rank distinction appears in the fiber Fourier analysis of
\cite[Section~4, especially Lemma~4.3]{StrombergssonVishe}.

The trivial representation of an amenable group is weakly
contained in its regular representation. Induction preserves
weak containment and takes the regular representation of a closed
subgroup to that of $G$. Thus every nonzero-frequency summand
in~\eqref{eq:markedtorusinduction} is tempered. See
\cite[Lemma 2.3 and Remark 2.4]{BenoistKobayashi}.
The whole orthogonal complement of the fiber average is tempered,
including its countably many frequency orbits. This gives uniform
control of their sum; there is no frequency-dependent implied
constant left to sum. Explicitly, for every $f\in C_c(G)$, all
summands satisfy $\|\pi_{[M]}(f)\|\leq\|\lambda_G(f)\|$;
their Hilbert direct sum satisfies the same inequality.

This simple stabilizer argument is enough for the present rate
estimate. It does not claim a complete irreducible spectral
decomposition. The corresponding finer decomposition for one
relative marked vector is developed in
\cite{AthreyaLagaceMollerRaum}; in particular, that paper's genuine
part contains no complementary series.

\subsubsection*{The base spectrum at level two}

The quotient $\Gamma(2)\backslash\HH$ is the thrice-punctured
sphere. The finite-area theorem of
\cite[Theorem 2]{OtalRosas} states that the discrete eigenvalue
$\lambda_{2g-2+n}$ exceeds $1/4$, counting $\lambda_0=0$ and
including residual as well as cuspidal eigenvalues. For $(g,n)=(0,3)$,
there is therefore no positive discrete eigenvalue at or below
$1/4$. The continuous spectrum starts at $1/4$.
Since $-I\in\Gamma(2)$, the base representation factors through
$\operatorname{PSL}_2(\RR)$. All non-tempered nontrivial representations of
$\operatorname{PSL}_2(\RR)$ in this setting are spherical
complementary-series representations and would produce a
Laplacian eigenvalue in $(0,1/4)$; this is the $d=1$ specialization of
the unitary-dual description in~\cite[Section~1]{EdwardsOhMixing}.
It follows that
$L^2_0(G/\Gamma(2))$ is tempered. The same conclusion holds for
$\Gamma=\operatorname{SL}_2(\ZZ)$, by pullback from that quotient
to the level-two cover.

Combining the base and nonzero-frequency parts proves that
\begin{equation}
 L^2_0(Y_2)\quad\hbox{is a tempered $G$-representation}
 \label{eq:markedtorusalltempered}
\end{equation}
for either of these two choices of $\Gamma$. The statement does
not automatically hold for a further finite-index subgroup.

\subsection{A global correlation bound with an explicit constant}

Let $K=\operatorname{SO}(2)$, with angular generator $\Theta$,
normalized so that the $n$th character has derivative $in$.
Use the norm
\begin{equation}
 \mathcal S_1(F)^2=\|F\|_2^2+\|\Theta F\|_2^2.
 \label{eq:markedtorusKnorm}
\end{equation}
The angle is the ordinary matrix rotation angle: on the orientation
cover, $\omega$ changes to $e^{i\theta}\omega$, so the genus-two
quadratic differential changes to $e^{2i\theta}q_X$. Thus this
normalization is twice differentiation in the phase of $q_X$.
Write $F^\circ=F-\mu(F)$ and $G^\circ=G-\mu(G)$.
The Cowling--Haagerup--Howe bound for a tempered representation is
\[
 |\langle\pi(g)v,w\rangle|
 \leq\Xi(g)\|v\|\|w\|
       \sqrt{\dim Kv\,\dim Kw}
\]
for $K$-finite vectors; see
\cite[Proposition 2.6]{BenoistKobayashi} and
\cite{CowlingHaagerupHowe}. For single $K$-types the two
dimension factors equal one. Decomposing into $K$-types and
using Cauchy--Schwarz gives
\[
 \sum_{n\in\ZZ}\|F_n\|_2
 \leq\left(\sum_{n\in\ZZ}(1+n^2)^{-1}\right)^{1/2}
       \mathcal S_1(F)
 =\sqrt{\pi\coth\pi}\,\mathcal S_1(F).
\]
Here the subscript denotes $K$-type, distinct from the matrix
frequency in~\eqref{eq:markedtorusfourier}.

For $e_1=(1,0)^T$, the spherical principal-series realization gives
$\Xi(a_t)=\int_K|a_tke_1|^{-1}\,\dd k$, with Haar probability on $K$.
Using $u=\tan\theta$ on a quadrant gives, in our time normalization,
the elementary integral
\begin{equation}
 \Xi(a_t)=\frac2\pi e^{-t}
   \int_0^\infty
    \frac{\dd u}{\sqrt{(1+u^2)(1+e^{-4t}u^2)}}
 \leq\frac4\pi(1+t)e^{-t},\qquad t\geq0.
 \label{eq:markedtorusXi}
\end{equation}
For the inequality, split the integral at $u=1$ and $u=e^{2t}$.
The three contributions are at most $1$, $2t$, and $1$.
\begin{theorem}[An explicit marked-torus correlation bound]
\label{thm:markedtorusmixing}
For Haar probability measure on $Y_2$ and all $F,G$ in the domain
of the rotation generator, the following inequality holds for $t\geq0$:
\begin{equation}
 \begin{split}
 \left|\int_{Y_2}\overline F\,(G\circ a_t)\,\dd\mu
       -\overline{\mu(F)}\mu(G)\right|
 \leq4\coth\pi\,(1+t)e^{-t}
       \mathcal S_1(F^\circ)\mathcal S_1(G^\circ).
 \end{split}
 \label{eq:markedtorusexplicitmix}
\end{equation}
The same estimate holds on $\mathcal Q^{\mathrm{ns},1}_2(2,2)$
with its probability-normalized Masur--Veech measure and the induced
rotation norm.
\end{theorem}
\begin{proof}
The preceding $K$-type estimate and~\eqref{eq:markedtorusXi} apply
to the tempered representation~\eqref{eq:markedtorusalltempered}.
Remove the constant components to obtain the displayed covariance.
The geometric statement follows by pulling observables back to the
finite model of Section~\ref{subsec:genus2homogeneous}; probability
normalization preserves the norms and the correlation.
\end{proof}
Centering the functions does not increase $\mathcal S_1$.
The norm needs only one rotation derivative; neither derivatives in
the marked-point coordinates nor compact support are required.
This is global on the finite-volume affine quotient. No compact
cutoff or unquantified cusp-tail estimate enters this inequality.
The polynomial factor allows every pure exponential rate
$e^{-(1-\varepsilon)t}$, but not a uniform constant times $e^{-t}$
from this argument. More explicitly, for $0<\varepsilon\leq1$ one
may replace $4\coth\pi(1+t)e^{-t}$ by
$4\coth\pi\,\varepsilon^{-1}e^{-(1-\varepsilon)t}$.
For the convention $a_t=\diag(e^{t/2},e^{-t/2})$, replace $t$
by $t/2$ throughout.

\subsection{An exact finite lattice sum with two endpoint weights}

The following unfolding supplies a concrete counting statistic.
For $\varphi,\psi\in C_c^\infty(\mathsf A)$ let
\[
 P\varphi(h\Lambda)=\sum_{\lambda\in\Lambda}\varphi(h\lambda),
 \qquad F=P\varphi,\quad G=P\psi.
\]
Normalize Haar measure $\dd h$ so that $\mathsf A/\Lambda$ has volume one.
For fixed $t$ define the overlap weight
\begin{equation}
 W_t^{\varphi,\psi}(\lambda)
 =\int_{\mathsf A}\overline{\varphi(h)}\,
                   \psi(a_th\lambda)\,\dd h.
 \label{eq:markedtorusweight}
\end{equation}
Only finitely many $\lambda$ contribute, since any such lattice
element lies in the compact set
$(\operatorname{supp}\varphi)^{-1}a_{-t}
\operatorname{supp}\psi$. In coordinates
$h=(g,v)$ and $\lambda=(\gamma,m)$, the weight is explicitly
\begin{equation}
 W_t^{\varphi,\psi}(\gamma,m)
 =\int_G\int_{\operatorname{Mat}_{2\times2}(\RR)}
    \overline{\varphi(g,v)}\,
    \psi\bigl(a_tg\gamma,a_t(v+gm)\bigr)\,\dd v\,\dd g.
 \label{eq:markedtorusweightcoordinates}
\end{equation}
Unfolding the periodization of the first observable gives exactly
\begin{equation}
 \sum_{\lambda\in\Lambda}W_t^{\varphi,\psi}(\lambda)
   =\int_{Y_2}\overline F\,(G\circ a_t)\,\dd\mu.
 \label{eq:markedtorusexactcount}
\end{equation}
The weights are nonnegative when $\varphi$ and $\psi$ are
nonnegative. For arbitrary observables they are signed or complex.
Thus~\eqref{eq:markedtorusexactcount} is a finite weighted return
sum, not an assertion about an unweighted sharp orbit count.

The unstable Jacobian exponent is $4$: conjugation by $a_t$
expands the upper-unipotent direction in $\mathfrak{sl}_2$ by
$e^{2t}$ and the two expanding translation directions by $e^t$
each. Accordingly define the normalized return statistic and its
cumulative version, for $T\geq1$, by
\[
 \mathcal N_{\varphi,\psi}(t)
   =e^{4t}\sum_{\lambda\in\Lambda}
                     W_t^{\varphi,\psi}(\lambda),\qquad
 \mathcal A_{\varphi,\psi}(T)
   =\int_1^T\mathcal N_{\varphi,\psi}(t)\,\dd t.
\]
This normalization uses the affine model's unstable Jacobian. It
does not use the polar Jacobian of full Teichm\"uller space.
Since $\mu(P\varphi)=\int_{\mathsf A}\varphi\,\dd h$,
\eqref{eq:markedtorusexplicitmix} gives the explicit remainder
\begin{equation}
 \begin{split}
 \left|\mathcal A_{\varphi,\psi}(T)
    -\overline{\mu(F)}\mu(G)\frac{e^{4T}-e^4}{4}\right|
 \leq\frac43\coth\pi\,(1+T)e^{3T}
       \mathcal S_1(F^\circ)\mathcal S_1(G^\circ).
 \end{split}
 \label{eq:markedtoruscountremainder}
\end{equation}
Equivalently, with $R=e^T$, this is an $O(R^3\log R)$
remainder against an $R^4$ main term, with a displayed norm bound.
Here $R=e^T$ is an exponential-time parameter; identifying it with
hyperbolic curve length, intersection number, or square-tiled area
would require a separate comparison. The estimate follows simply from
$\int_1^T(1+t)e^{3t}\dd t\leq\tfrac13(1+T)e^{3T}$.

These norms can also be bounded using the original compactly
supported weights. If $K_0\subset \mathsf A$ contains their supports and
\[
 M_{K_0}=\#(\Lambda\cap K_0^{-1}K_0),
\]
then Cauchy--Schwarz in the locally finite periodization gives
\[
 \mathcal S_1(P\varphi)
  \leq\sqrt{M_{K_0}}
       \bigl(\|\varphi\|_{L^2(\mathsf A)}^2+
              \|\Theta\varphi\|_{L^2(\mathsf A)}^2\bigr)^{1/2},
\]
and similarly for $\psi$. All quantities in this bound are finite
for the chosen compact supports.

Define the flow generator by
$V\psi(h)=\left.\frac{\dd}{\dd s}\right|_{s=0}\psi(a_sh)$.
Periodization commutes with this generator:
$V(P\psi)=P(V\psi)$. Therefore inserting the terminal weight
$(4+V)\psi$ gives the exact identity
\begin{equation}
 \begin{split}
 \mathcal A_{\varphi,(4+V)\psi}(T)
  &=e^{4T}\int_{Y_2}\overline F\,(G\circ a_T)\,\dd\mu\\
  &\quad-e^4\int_{Y_2}\overline F\,(G\circ a_1)\,\dd\mu.
 \end{split}
 \label{eq:markedtorusderivativeweight}
\end{equation}
There is no convolution operator to commute past $V$, and no
unjustified differentiation of a counting asymptotic.
The quantitative norm required for this terminal test is
$\mathcal S_1((4+V)G)\leq4\mathcal S_1(G)+\mathcal S_1(VG)$.
It includes $\|\Theta VG\|_2$; the generators $\Theta$ and $V$
do not commute, and this mixed derivative is not omitted.


\subsection{What the marked-torus calculation settles}
\label{subsec:genus2torusscope}

The compact polyhedron encodes the genus-two hyperbolic base, whereas the
marked flat torus encodes the nonsquare $(2,2)$ differential and its flow.
They use different intrinsic metrics and are linked by the same six branch
values. The affine model gives an exact Fourier decomposition, an explicit
global correlation bound, and the finite weighted return sum with
remainder~\eqref{eq:markedtoruscountremainder}. These are conclusions of the
calculation, not hypotheses imported from general Teichm\"uller mixing.
The inputs are classical: amenability and induction, the
Cowling--Haagerup--Howe bound, the elementary $SL_2$ spherical function,
and the absence of positive small spectrum on the thrice-punctured sphere.
Related marked-torus harmonic analysis is already developed in
\cite{StrombergssonVishe,AthreyaLagaceMollerRaum}; no novelty claim for that
machinery is intended.

The derivative identity~\eqref{eq:markedtorusderivativeweight} also makes
the counting converse exact in this model. After subtracting the product
of means, a bound for the terminal test $(4+V)G$ recovers the correlation
at $T$ with the same exponential saving. In the proof here the direction
is Fourier/representation theory to correlations and hence to counts.
The identity identifies the reverse implication as well, without making
the proof circular.

The restriction to the chosen stratum remains substantial. Although it
meets every genus-two conformal base, its unit-area dimension is seven;
the principal unit stratum has dimension eleven. The affine expanding
Jacobian is $e^{4t}$, whereas the full genus-two Teichm\"uller polar kernel
in Section~\ref{sec:pairconverse} has exponent $6$. The two kernels cannot
be substituted for one another. Section~\ref{sec:coefficientobservables}
constructs selected primitive-arc observables from the coefficient data
with controlled geometric norms. A comparison with the full orbit kernel,
and uniform stability of the nonlinear coefficient extraction, remain
open. In particular, the $O(R^3\log R)$ estimate above is not an error
term for the unweighted count of connected simple closed geodesics.
A comparison with that count, or an extension to the principal stratum
$\mathcal Q^1_2(1^4)$, is a further problem.



